\documentclass[../NDCNNAI_main.tex]{subfiles}
\graphicspath{{\subfix{../fig/}}}

\begin{document}

	\subsection{Введение в функциональные данные: функции как объекты}
	
	В \emph{функциональном анализе данных} (Functional Data Analysis, FDA) каждый объект (индивид, эксперимент) описывается не вектором из $\mathbb{R}^n$, а одной или несколькими функциями $g\colon Z \to \mathbb{R}$, где $Z \subset \mathbb{R}^u$ --- компакт (например, временной интервал $Z = [-T, 0]$).
	
	\textbf{Зачем это нужно для ИНС?}
	\begin{itemize}
		\item Естественный учёт зависимостей между измерениями (например, временная корреляция в сигналах).
		\item Работа с неравномерной или разреженной выборкой, редкими измерениями: функция (сигнал) восстанавливается через сглаживание (сплайны, ядра, вейвлеты).
		\item Возможность применять классические задачи машинного обучения (регрессия, классификация, кластеризация) к функциональным предикторам.
	\end{itemize}
	
	%------------------------------------------------------
	
	\subsection{Функциональные нейроны, архитектура FMLP и связь с обычными MLP}
	
	Обычный нейрон в MLP вычисляет $N(x) = T(w \cdot x + b)$, где $w \cdot x$ --- линейная форма (скалярное произведение) на $\mathbb{R}^n$. Для функционального входа $g \in L^p(Z, \mu)$ (или $C(Z)$) линейную форму естественно задать через интеграл.
	
	\begin{definition}[Функциональный нейрон]
		Пусть заданы функция активации $T\colon \mathbb{R} \to \mathbb{R}$, функциональные веса $f_\ell$ (измеримые функции, такие что $f_\ell g_\ell \in L^1(\mu_\ell)$) и смещение $b \in \mathbb{R}$. \emph{Функциональный нейрон} с $s$ входами $g_1,\dots,g_s$ определяется как
		\[
		N(g_1,\dots,g_s) = T\left( b + \sum_{\ell=1}^s \int_{Z_\ell} f_\ell(x) \, g_\ell(x) \, d\mu_\ell(x) \right).
		\]
	\end{definition}
	Проблема этого определения в том, что тут нечего обучать, веса $f_\ell$ заранее заданы, поэтому вводится параметризация.
	\begin{definition}[Параметрический функциональный нейрон]
		Пусть заданы параметрические семейства функций $F_\ell(w_\ell, \cdot)\colon Z_\ell \to \mathbb{R}$, где $w_\ell \in W_\ell \subset \mathbb{R}^{v_\ell}$ --- параметры, и $F_\ell(w_\ell, \cdot) \in L^{q_\ell}(\mu_\ell)$ (где $q_\ell$ --- сопряжённый показатель к $p_\ell$). Тогда \emph{параметрический функциональный нейрон} имеет вид:
		\[
		N(g_1,\dots,g_s) = T\left( b + \sum_{\ell=1}^s \int_{Z_\ell} F_\ell(w_\ell, x) \, g_\ell(x) \, d\mu_\ell(x) \right).
		\]
	\end{definition}
	
	Функциональные нейроны используются только в первом слое, последующие слои --- обычные числовые MLP. Таким образом, FMLP --- это композиция линейного интегрального оператора (извлекающего конечное число признаков из функции) и стандартной нейронной сети.
	
	\begin{definition}[Однослойный FMLP с одним функциональным входом]
		Для функционального входа $g$ и $k$ скрытых функциональных нейронов:
		\[
		H(g) = \sum_{i=1}^{k} a_i \, T\left( b_i + \int_{Z_i} F_i(w_i, x) g(x) d\mu(x) \right),
		\]
		где $a_i, b_i \in \mathbb{R}$, $w_i \in W_i$ --- параметры.
	\end{definition}
	
	\textbf{Связь с обычными MLP:}
	\begin{itemize}
		\item Если вместо функции $g$ подставить вектор $\mathbf{x} \in \mathbb{R}^n$, а меру $\mu$ взять дискретную с массами в $x_j$, интеграл превращается в сумму:
		\[
		\int_Z F(w, x) g(x) d\mu(x) \approx \sum_{j=1}^n F(w, x_j) g(x_j),
		\]
		что соответствует линейной форме $w \cdot g$ с $w_j = F(w, x_j)$.
		
		\item Таким образом, обычный MLP --- частный случай FMLP, когда функциональный вход задан только в конечном числе точек.
		
	\end{itemize}
	
	%---------------------------------------------------------

\subsection{Основная УАТ для FMLP (на $L^p$, $L^1$, произведениях функциональных пространств)}

Пусть \( C(A, B) \) обозначает множество непрерывных отображений из топологического пространства \( A \) в топологическое пространство \( B \). Для компактного подмножества \( K \subset X \) (где \( X \) --- топологическое векторное пространство) мы снабжаем пространство \( C(K, \mathbb{R}) \) \textbf{равномерной (чебышёвской) нормой}:
\[
\rho_K(f, g) \coloneq \sup_{x \in K} |f(x) - g(x)|.
\]

Если \( A = \mathbb{R}^n \), а \( M^n \) обозначает (борелевские) измеримые отображения \( \mathbb{R}^n \to \mathbb{R} \), то стандартной метрикой, порождающей сходимость на компактах, является
\[
d_C(f, g) = \sum_{m=1}^{\infty} 2^{-m} \min \left\{ \sup_{\|x\| \leq m} |f(x) - g(x)|, 1 \right\}.
\]

Для корректной формулировки теорем о плотности множества функций (например, нейронных сетей) в более широком пространстве вводится следующее техническое определение.

\begin{definition}[Внутренняя/внешняя плотность (по Стинчкомбу, 1999)]
Пусть \( (\mathcal{X}, d) \) --- метрическое пространство, \( C \subset \mathcal{X} \) и \( S \subset \mathcal{X} \).
\begin{itemize}
    \item Множество \( S \) называется \textbf{\( d \)-внешне плотным} в \( C \), если \( d \)-замыкание множества \( S \) содержит \( C \), то есть $C \subset \overline{S}^{\,d}$.
    \item Множество \( S \) называется \textbf{\( d \)-внутренне плотным} в \( C \), если \( d \)-замыкание множества \( S \cap C \) содержит \( C \), то есть $C \subset \overline{S \cap C}^{\,d}$.
\end{itemize}
Если \( C = \mathcal{X} \), эти понятия совпадают, и мы просто говорим, что \( S \) \textbf{\( d \)-плотно} в \( \mathcal{X} \).
\end{definition}

Пусть \( X \) --- топологическое векторное пространство, а \( X^* \) --- его (непрерывное) сопряжённое пространство. Пусть \( T\colon \mathbb{R} \to \mathbb{R} \) --- фиксированная функция активации.

\begin{definition}[Классический однослойный перцептрон]
Для размерности \( n \in \mathbb{N} \) определим класс функций
\[
S_T^n \coloneq \left\{ h(x) = \sum_{i=1}^{p} \beta_i \, T(w_i \cdot x + b_i) \;:\; p \in \mathbb{N},\; \beta_i \in \mathbb{R},\; (w_i, b_i) \in \mathbb{R}^n \times \mathbb{R} \right\}.
\]
\end{definition}

\begin{definition}[Обобщённый перцептрон на \( X \) (в стиле Стинчкомба)]
Пусть \( A \subset X^* \) --- некоторое множество непрерывных линейных функционалов на \( X \). Определим класс функций
\[
S_T^X(A) \coloneq \left\{ h(x) = \sum_{i=1}^{m} \beta_i \, T\bigl( \ell_i(x) + b_i \bigr) \;:\; m \in \mathbb{N},\; \beta_i, b_i \in \mathbb{R},\; \ell_i \in A \right\}.
\]
Здесь \( \ell_i\colon X \to \mathbb{R} \) --- линейные функционалы из множества \( A \), заменяющие скалярные произведения \( w_i \cdot x \).
\end{definition}

Пусть $\mu$ --- конечная борелевская мера на $\mathbb{R}^n$, и $1 < p < \infty$, а $q$ --- сопряжённый показатель ($\frac{1}{p} + \frac{1}{q} = 1$). Рассмотрим пространство $X = L^p(\mu)$ и пусть $V \subset L^q(\mu)$ --- плотное подмножество (в норме $L^q$). Определим множество интегральных линейных функционалов:
\[
A_V \coloneq \left\{ \ell_g \colon f \mapsto \int g(x)f(x) \, d\mu(x) \; \middle| \; g \in V \right\} \subset X^*.
\]

\begin{lemma}[О слабой-$\ast$ плотности интегральных форм]
Множество $A_V$ является \textbf{слабо-$\ast$ плотным} в пространстве $X^*$, которое изометрично $L^q(\mu)$ (по теореме Рисса о представлении).
\end{lemma}

\begin{proof}[Идея доказательства]
По теореме Рисса, любой линейный непрерывный функционал $\ell \in X^*$ имеет единственное представление в виде $\ell(f) = \int g_\ell(x)f(x) d\mu(x)$ для некоторой $g_\ell \in L^q(\mu)$. Так как $V$ плотно в $L^q(\mu)$, для любого $g_\ell$ найдётся последовательность $\{g_n\} \subset V$, сходящаяся к $g_\ell$ в норме $L^q$. Тогда соответствующие функционалы $\ell_{g_n}$ сходятся к $\ell$ в слабой-$\ast$ топологии.
\end{proof}

\textbf{Случай произведения пространств.} Для $X = \prod_{j=1}^s L^{p_j}(\mu_j)$ и плотных подмножеств $V_j \subset L^{q_j}(\mu_j)$ ($q_j$ --- сопряжённые к $p_j$ показатели) определим:
\[
A_V \coloneq \left\{ \ell_g(f_1, \ldots, f_s) = \sum_{j=1}^s \int g_j(x)f_j(x) \, d\mu_j(x) : g_j \in V_j \right\} \subset X^*.
\]
Тогда $A_V$ является слабо-$\ast$ плотным в $X^* \cong \prod_{j=1}^s L^{q_j}(\mu_j)$.

Далее мы обсудим общую схему доказательства плотности в топологических векторных пространствах, предложенную Стинчкомбом.

\begin{theorem}[Общая схема обоснования плотности (по Стинчкомбу, 1999)]
\label{thm:stinchcombe}
Пусть $X$ --- локально выпуклое топологическое векторное пространство и $A \subset X^*$. Предположим, что:
\begin{enumerate}
    \item $A$ является \textbf{слабо-$\ast$ плотным} в $X^*$ (т.е. его замыкание в топологии $\sigma(X^*, X)$ совпадает с $X^*$);
    \item Функция активации $T\colon \mathbb{R} \to \mathbb{R}$ \textbf{не является полиномом} (эквивалентно, является дискриминаторной в смысле Цыбенко/Лешно).
\end{enumerate}
Тогда для любого компактного подмножества $K \subset X$ множество $S_T^X(A)$ является \textbf{равномерно плотным} в пространстве непрерывных функций $C(K, \mathbb{R})$:
\[
\overline{S_T^X(A) \big|_K}^{\,\rho_K} = C(K, \mathbb{R}),
\]
где $\rho_K$ --- равномерная норма на $K$.
\end{theorem}

\begin{proof}
Доказательство основано на \textbf{теореме 5.1} и \textbf{следствии 5.1.3} из статьи Stinchcombe (1999).

\textbf{1. Переход к одномерному случаю.} По условию (2) функция активации $T$ не является полиномом. Согласно классическому результату Лешно--Лина--Пинкуса--Шокена (1993), множество 
    \[
    S_T^1 = \mathrm{span}\{ T(w \cdot x + b) \mid w, b \in \mathbb{R} \}
    \]
    является плотным в $C(\mathbb{R})$ в топологии равномерной сходимости на компактах ($d_C$-плотным). В частности, $S_T^1$ $d_C$-внешне плотно в $C^1$.

\textbf{2. Использование слабой-$\ast$ плотности $A$.} 
    Рассмотрим пространство $X^*$ с топологией $\sigma(X^*, X)$. По условию (1) $A$ слабо-$\ast$ плотно в $X^*$. 
    
    Для любого компакта $K \subset X$ рассмотрим семейство функций 
    \[
    \mathcal{A} = \{ \ell(\cdot) + b \mid \ell \in A, b \in \mathbb{R} \}.
    \]
    Это векторное подпространство в $C(K)$, так как оно замкнуто относительно сложения и умножения на скаляр.

\textbf{3. Проверка условий теоремы 5.1(b).} 
    \begin{itemize}
        \item \textbf{Постоянные функции:} выбирая $\ell = 0 \in \overline{A}$ (слабое-$\ast$ замыкание), получаем постоянные функции.
        \item \textbf{Разделение точек:} для любых $x \neq y$ из $K$ найдётся $\ell \in A$ такой, что $\ell(x) \neq \ell(y)$ (по теореме Хана-Банаха и слабой-$\ast$ плотности $A$). Добавляя сдвиг $b$, можно добиться разделения.
    \end{itemize}
    Таким образом, $\mathcal{A}$ удовлетворяет условиям теоремы 5.1(b).

\textbf{4. Применение теоремы 5.1.} 
    По теореме 5.1(b), если $S_T^1$ является $d_C$-внешне плотным в $C^1$, то $\mathrm{span}(T \circ \mathcal{A})$ будет $\rho_K$-внешне плотным в $C(K)$. Но
    \[
    \mathrm{span}(T \circ \mathcal{A}) = S_T^X(A),
    \]
    что и означает равномерную плотность на компакте $K$.

\textbf{4. Завершение доказательства.} 
    Поскольку $K$ произвольный компакт в $X$, получаем требуемое:
    \[
    \overline{S_T^X(A) \big|_K}^{\,\rho_K} = C(K, \mathbb{R}).
    \]

\end{proof}

\begin{theorem}[Функциональный аналог УАТ в \( L^p \)]
\label{thm:functional-UAT-Lp}
Пусть \( 1 < p < \infty \), \( X = L^p(\mu) \) с конечной борелевской мерой \(\mu\), и \( V \subset L^q(\mu) \) --- плотное подмножество. Предположим, что \( T: \mathbb{R} \to \mathbb{R} \) не является полиномом (так что \( S_T^1 \) \( d_C \)-плотно в \( C(K_1, \mathbb{R}) \) для любого компакта \( K_1 \subset \mathbb{R} \)). Тогда для любого компакта \( K \subset X \)

\[
\overline{S_T^X(A_V)|_K}^{\,\rho_K} = C(K, \mathbb{R}), \qquad 
A_V = \left\{ \ell_g(f) = \int g f \, d\mu : g \in V \right\}.
\]
\end{theorem}

\begin{proof}
    \textbf{1. Отождествление двойственного пространства.} По теореме Рисса, \( X^* \simeq L^q(\mu) \). Множество \( A_V \) состоит из функционалов вида \( \ell_g(f) = \langle g, f \rangle \), где \( g \in V \subset L^q \).

    \textbf{2. Слабая-$\ast$ плотность \( A_V \).} Согласно лемме, \( A_V \) слабо-$\ast$ плотно в \( X^* \).

    \textbf{3. Применение теоремы Стинчкомба.} Проверим условия предыдущей \hyperref[thm:stinchcombe]{теоремы}:
    \begin{itemize}
        \item \( X = L^p(\mu) \) --- локально выпуклое топологическое векторное пространство (нормированное, рефлексивное);
        \item \( A_V \) слабо-$\ast$ плотно в \( X^* \) согласно шагу 2;
        \item \( T \) не является полиномом --- условие активации выполнено.
    \end{itemize}
    Тогда по теореме Стинчкомба множество 
    \[
    S_T^X(A_V) = \mathrm{span}\{ T(\ell_g(\cdot) + b) : \ell_g \in A_V, \, b \in \mathbb{R} \}
    \]
    равномерно плотно в \( C(K, \mathbb{R}) \) для любого компакта \( K \subset X \).

    \textbf{4. Завершение доказательства.} Плотность \( S_T^X(A_V) \) в \( C(K) \) означает, что любую непрерывную функцию на \( K \) можно равномерно приблизить функциями вида
    \[
    f \mapsto \sum_{k=1}^N \beta_k \, T\!\left( \int g_k(x) f(x) \, d\mu(x) + b_k \right), \quad g_k \in V, \; b_k, \beta_k \in \mathbb{R}.
    \]
    Это и есть утверждение теоремы.
\end{proof}
Рассмотрим теперь случай $L^1$, выделяем его отдельно из-за бесконечности сопряжённого показателя.
\begin{theorem}[Функциональная УАТ в \( L^1 \) с компактным носителем меры]
\label{thm:functional-UAT-L1}
Пусть \( X = L^1(\mu) \), где \(\mu\) --- конечная борелевская мера с компактным носителем. 
Пусть \( V \subset L^\infty(\mu) \) слабо-$\ast$ плотно в \( L^\infty(\mu) \) в топологии \( \sigma(L^\infty, L^1) \), и пусть 
\[
A_V = \left\{ \ell_g(f) = \int g f \, d\mu : g \in V \right\}.
\]
Если \( T\colon \mathbb{R} \to \mathbb{R} \) не является полиномом, то для любого компакта \( K \subset X \)
\[
\overline{S_T^X(A_V)|_K}^{\,\rho_K} = C(K, \mathbb{R}).
\]
\end{theorem}

\begin{proof}
    \textbf{1. Двойственное пространство.} Для \( L^1(\mu) \) с конечной мерой \( X^* = L^\infty(\mu) \). Топология на \( L^\infty \), согласованная с двойственностью \(\langle L^\infty, L^1 \rangle\), --- это слабая-$\ast$ топология \( \sigma(L^\infty, L^1) \).

    \textbf{2. Слабая-$\ast$ плотность \( A_V \).} По условию \( V \) слабо-$\ast$ плотно в \( L^\infty(\mu) \), следовательно, \( A_V \) слабо-$\ast$ плотно в \( X^* \).

    \textbf{3. Компактность носителя меры.} Требование компактности носителя \(\mu\) обеспечивает выполнение условий теоремы Стинчкомба: 
    \begin{itemize}
        \item Пространство \( L^1(\mu) \) с такой мерой сепарабельно и обладает свойством аппроксимации.
        \item Компактность носителя позволяет использовать стандартные рассуждения типа Стоуна-Вейерштрасса для алгебр функций, порождённых \( T \circ \ell_g \).
    \end{itemize}

    \textbf{4. Применение общей теоремы.} Условия теоремы Стинчкомба выполнены:
    \begin{itemize}
        \item \( X = L^1(\mu) \) --- локально выпуклое топологическое векторное пространство;
        \item \( A_V \) слабо-$\ast$ плотно в \( X^* \);
        \item \( T \) --- не полиномиальная активация.
    \end{itemize}
    Тогда по \hyperref[thm:stinchcombe]{теореме} множество \( S_T^X(A_V) \) равномерно плотно в \( C(K, \mathbb{R}) \) для любого компакта \( K \subset X \).
\end{proof}
Теперь мы готовы доказать основной результат про универсальную аппроксимацию для функциональных многослойных перцептронов.
\begin{theorem}[УАТ для FMLP на произведении пространств]
\label{thm:main-UAT-FMLP}
Пусть \( X = \prod_{j=1}^s L^{p_j}(\mu_j) \), где \(\mu_j\) --- конечные борелевские меры, и \( 1 < p_j < \infty \) (или \( p_j = 1 \) с условиями компактного носителя / слабой-$\ast$ плотности из теоремы \ref{thm:functional-UAT-L1}).  
Для каждого \( j \) пусть \( V_j \subset L^{q_j}(\mu_j) \) плотно (или \( V_j \subset L^\infty \) слабо-$\ast$ плотно для \( p_j = 1 \)).  
Определим множество линейных форм:

\[
A_V = \left\{ \ell_g(f_1,\dots,f_s) = \sum_{j=1}^s \int g_j(x) f_j(x) \, d\mu_j(x) : g_j \in V_j \right\} \subset X^*.
\]

Пусть \( T\colon \mathbb{R} \to \mathbb{R} \) --- не полиномиальная функция (например, \(\th\), сигмоида, \(\mathrm{ReLU}^\alpha\) с нецелым \(\alpha\) и т.д.).  
Тогда для любого компакта \( K \subset X \)

\[
\overline{S_T^X(A_V)|_K}^{\,\rho_K} = C(K, \mathbb{R}).
\]

Иными словами, FMLP с одним скрытым слоем и интегральными линейными функционалами в качестве входов обладают свойством УА на компактных подмножествах \( X \).
\end{theorem}

\begin{proof}
    \textbf{1. Слабая-$\ast$ плотность \( A_V \).}  
    Из леммы о произведении пространств следует, что \( A_V \) слабо-$\ast$ плотно в \( X^* = \prod_{j=1}^s L^{q_j}(\mu_j) \) (или \( L^\infty \) при \( p_j = 1 \)).

    \textbf{2. Проверка условий теоремы Стинчкомба.}  
    \begin{itemize}
        \item \( X \) --- локально выпуклое топологическое векторное пространство (как произведение банаховых пространств).
        \item \( A_V \) слабо-$\ast$ плотно в \( X^* \) по шагу 1.
        \item \( T \) не является полиномом --- условие активации выполнено.
    \end{itemize}

    \textbf{2. Применение общей теоремы.}  
    По теореме Стинчкомба множество 
    \[
    S_T^X(A_V) = \mathrm{span}\{ T(\ell_g(\cdot) + b) : \ell_g \in A_V, \; b \in \mathbb{R} \}
    \]
    является равномерно плотным в \( C(K, \mathbb{R}) \) для любого компакта \( K \subset X \).

    \textbf{3. Интерпретация как FMLP.}  
    Элементы \( S_T^X(A_V) \) имеют вид
    \[
    f \mapsto \sum_{k=1}^N \beta_k \, T\!\left( \sum_{j=1}^s \int g_{jk}(x) f_j(x) \, d\mu_j(x) + b_k \right),
    \]
    что соответствует однослойной нейросети, где входы --- линейные интегральные функционалы от компонент \( f = (f_1,\dots,f_s) \).
\end{proof}

\subsection{Сопоставление с результатами лекций 1--4 (классические УАТ, GTN, операторные сети)}
\label{sec:big-picture}

Структура курса естественным образом приводит к функциональной теореме универсальной аппроксимации (УАТ) через следующую последовательность шагов:

\begin{itemize}
    \item \textbf{Лекции 1--2:} Классическая УАТ на \(\mathbb{R}^n\) (Цыбенко/Лешно) даёт \textbf{скалярную универсальность} функции активации \(T\):
    \[
    \overline{\mathrm{span}\{T(w \cdot x + b)\}}^{C(K)} = C(K, \mathbb{R}), \quad K \subset \mathbb{R}^n \text{ компакт}.
    \]

    \item \textbf{Лекции 3--4:} Операторный взгляд (GTN / выбор точек и работа с значениями в них) показывает, как линейные функционалы извлекают конечное количество <<состояние/признаков>> из сигналов:

    \item \textbf{Лекция 5 (текущая):} Выбираем \(A\) как множество интегральных линейных форм на \(L^p\) (теорема Рисса), убеждаемся, что \(A\) слабо-$\ast$ плотно в \(X^*\), и \textbf{переносим} скалярную УАТ на функциональные входы через схему Стинчкомба:
    \[
    \overline{\mathrm{span}\{T(\ell(\cdot) + b) : \ell \in A\}}^{C(K)} = C(K, \mathbb{R}), \quad K \subset X \text{ компакт}.
    \]
\end{itemize}

Таким образом, ключевая идея может быть записана в виде следующей импликации:

\begin{figure}
\begin{center}
    \begin{tikzpicture}[
        node distance=0.5cm,
        every node/.style={align=center},
        box/.style={rectangle, draw=black, thick, minimum width=3.5cm, minimum height=0.8cm, align=center},
    ]

    % Узлы
    \node[box] (A) {Скалярная\\экспрессивность $T$};
    \node[right=0.1cm of A] (plus) {$+$};
    \node[right=0.1cm of plus, box] (B) {Слабо-$\ast$ плотные\\линейные формы $A \subset X^*$};
    \node[right=0.1cm of B] (implies) {$\Longrightarrow$};
    \node[right=0.1cm of implies, box] (C) {FMLP универсальность\\на $C(K,\mathbb{R})$};

    \end{tikzpicture}
\end{center}
\caption{FMLP подход к задаче универсальной аппроксимации.}
\end{figure}

\end{document}